% Comparação dos resultados simulados
\begin{frame}{Validação no LTspice}
  \renewcommand{\arraystretch}{0.92}
  \begin{columns}[T,onlytextwidth]
    \begin{column}{0.43\textwidth}
      \centering
      \includegraphics[
        page=10,
        trim=100bp 330bp 100bp 320bp,
        clip,
        width=\linewidth,
        height=0.46\textheight,
        keepaspectratio
      ]{\detokenize{prism-uploads/Romário_Jonas_-_relatório_Pratica_02___Circuitos_Elétricos_2.pdf}}

      \scriptsize Circuito no LTspice.
    \end{column}
    \begin{column}{0.54\textwidth}
      \centering
      \scriptsize
      \begin{tabular}{rccc}
        \toprule
        $f$ (Hz) & Teórico & LTspice & Erro \\
        \midrule
        40    & 0,5947 & 0,5952 & 0,09\% \\
        100   & 1,5431 & 1,5480 & 0,31\% \\
        400   & 18,1956 & 22,9160 & 25,94\% \\
        480   & 49,9733 & 48,7146 & 2,52\% \\
        1100  & 3,7874 & 2,9902 & 21,05\% \\
        11000 & 0,3084 & 0,2897 & 6,07\% \\
        \bottomrule
      \end{tabular}

      \vspace{0.16cm}
      \textbf{Fase teórica e simulada}

      \vspace{0.06cm}
      \begin{tabular}{rcc}
        \toprule
        $f$ (Hz) & Teórica & LTspice \\
        \midrule
        40    & $-90,68^\circ$ & $-89^\circ$ \\
        480   & $178,13^\circ$ & $179^\circ$ \\
        11000 & $90,35^\circ$  & $92^\circ$ \\
        \bottomrule
      \end{tabular}
    \end{column}
  \end{columns}

  \vspace{0.06cm}
  \begin{alertblock}{Síntese}
    \centering
    A simulação reproduziu a ressonância e as tendências de magnitude e fase;
    os maiores desvios ocorreram nas regiões de transição.
  \end{alertblock}
\end{frame}